\begin{pastexam}[label={px:nchu-111-is-09}]{中興資工所乙組 111 學年度 第 20 題}
（問答）

{[6\%]} Please describe the greedy algorithm and find an example in which the greedy algorithm does not find
the optimal solution.
\end{pastexam}

\begin{pastsoln}
\textbf{貪心法是什麼}。 貪心法把問題的解看成一連串的選擇，每一步都依某個規則挑「眼前看起來最好」的
選項，選了就不再改，再去解剩下的子問題（資結書第~19.1.6~節）。 它不像動態規劃那樣比較所有選擇的結果，
所以通常很快，但只在問題有兩個性質時才保證最佳：
\begin{itemize}
\item \textbf{貪心選擇性質}：存在一個最佳解包含第一步的貪心選擇（通常用交換論證證明）。
\item \textbf{最佳子結構}：做了貪心選擇之後，剩下的問題的最佳解，和這個選擇合起來就是原問題的最佳解。
\end{itemize}
成功的例子：活動選擇（每次選最早結束的活動）、霍夫曼編碼、Kruskal 和 Prim 的最小生成樹、分數背包
（依單位價值由高到低拿）。

\textbf{貪心法失敗的例子：找零錢}。 硬幣面額 $\{1, 3, 4\}$，要湊出 $6$ 元，硬幣數越少越好。 貪心規則
「每次拿不超過剩餘金額的最大面額」：先拿 $4$，剩 $2$；再拿 $1$、$1$，共 $3$ 枚（$4 + 1 + 1$）。 但
$3 + 3$ 只要 $2$ 枚。 貪心法第一步拿 $4$ 就錯了：包含 $4$ 的湊法最少要 $3$ 枚，而最佳解不含 $4$，貪心
選擇性質不成立（資結書第~19.5~節）。 這個問題要用動態規劃：$c(x) = 1 + \min_{d \le x} c(x - d)$。

另一個常見的例子是 0/1 背包：容量 $50$，物品（重量, 價值）為 $(10, 60)$、$(20, 100)$、$(30, 120)$。
依單位價值 $6, 5, 4$ 由高到低拿，拿了前兩件後剩 $20$ 放不下第三件，價值 $160$；但拿第二、三件是
$220$。 物品不能切開時，單位價值最高的物品不一定在最佳解裡。
\end{pastsoln}
